\section{Source Pipeline and Run Provenance}
\label{app:pipeline}

\subsection{Calibration, revision, and source interpretation}
Removing eight development--test article overlaps leaves 2,651 eligible development questions.
A fixed hash selects 176 questions from 50 articles for the ranking pilot.
The initial remainder contains 2,475 questions from 679 articles.
Earlier source inspections identify 35 further articles for exclusion before their first validation outcomes.
Removing their 127 questions gives the 2,348-question template track.
The human track retains 830 positive questions from the published test collection.
Both tracks had inspected outcomes before the final source revision.
The final evaluation preserves their questions, ranking, budgets, reader settings, parser, and metrics.
The release identifies these runs as post-audit revision comparisons.

Initial source calibration covers 86 articles, including 85 with eligible development questions.
It combines 50 pilot articles with 39 previously identified source paths, with three overlaps.
Ranking selection uses only the 176-question pilot.
An independent feature audit changes query metadata and remaps source identifiers without changing selections.
Ranking accepts the question string alone.
The title inventory comes from the source corpus.
Every ranking retains every corpus unit exactly once.

A year maps to January 1 through December 31 on the Gregorian ordinal axis.
A month maps to its actual first and last dates.
A fully stated day gives a singleton.
Tenures keep separate uncertain start and end events.
The final adapter requires local subjects, preserves abbreviation boundaries, checks date ownership, and links uniquely resolved paragraph-level ends.
Recognized bounded states with unresolved endings remain fallback text.
No answer labels guide these repairs.
Twenty-seven source checks pass before the final evaluation.
Five inspected calendar-end cases receive explicit endpoints.
Six inspected duration cases remain fallback text because their endings cannot enter the independent-bound model.
These eleven cases guide development and do not estimate extraction accuracy.

All source words survive the modeled/fallback partition.
Independent checks find no overlapping units and no rejected pair endpoints.
The development and test pools contain 33,005 and 34,082 units, respectively.
Their modeled claim counts are 493 and 464, with 292 and 238 explicit ends.
Sixteen paragraph-level end links enter the two pools.
The development corpus contains six succession starts and three directed pairs.
The test corpus contains two succession starts and no pairs.
All three pairs occur in previously inspected calibration articles.
The Warburg example uses automatic extraction; its question period and ranking are illustrative.
Its selected narrative dates Bing's appointment to 1955, while another paragraph lists 1954.
The example compiles the stated narrative without resolving that cross-paragraph discrepancy.

\paragraph{Earlier source audit.}
The preceding frozen version was assessed on 60 development claims outside its calibration articles.
That model-assisted audit accepted 52 events, rejected five, and marked three uncertain.
Eleven accepted starts omitted terminations elsewhere in their paragraphs.
Twenty-seven of its 29 explicit-end records had accepted event grounding.
The final adapter addresses those findings through the five source-attachment checks described above.
The earlier judgments remain unchanged in the release.

\paragraph{Recorded source audit.}
After the adapter freeze, a fixed hash selects 60 claims from 55 articles.
The sample contains twenty development claims and forty test claims.
It excludes 277 documented or identifiable exposure articles.
The complete membership of a lost earlier audit is unavailable.
We cannot rule out prior exposure for every article.
Reviewers read complete source paragraphs without benchmark questions, labels, retrieval scores, or reader outputs.
They accept 54 core interpretations, reject four, and mark two uncertain.
All 408 checked spans match the source characters exactly.
Twenty-six of 27 supplied ends are accepted; one is rejected.
Core correctness and lifecycle completeness receive separate judgments.
Eight accepted events omit explicit paragraph-level ends; five further endings have uncertain applicability.
These model-assisted judgments are neither human annotations nor extraction-recall estimates.
The samples differ, so the earlier 52/60 and current 54/60 are not paired improvement estimates.
No judgment changes the final frozen adapter.

\subsection{Reader protocol}
\label{app:reader}
The current model is Qwen2.5-3B-Instruct, Q4\_K\_M, at revision \texttt{7dabda4d}.
The release records the full revision and weight checksum.
The runtime is llama.cpp b11435 at commit \texttt{43fe9c6}.
The model uses the Qwen Research License; the runtime uses MIT.
Generation weights are excluded from the archive.
Pinned download scripts restore the official artifacts.

Prompts contain the visible question, original excerpts, source titles, date bounds, and temporal status.
Modeled excerpts also carry source-derived subject and relation labels.
Those labels preserve role context for short succession clauses.
No extracted answer-value field is added to the prompt metadata.
Temperature is zero, seed is 1729, and the output limit is 96 tokens.
Prompt caching and JSON grammar enforcement are disabled.
A tokenizer check verifies that every request fits the configured context.
Resource-only amendments precede completed question-answer outputs and preserve the original freezes and logs.
The final runtime uses four threads, a 2,048-token context, and polling disabled.
Its batch and microbatch sizes are both 512.

The parser accepts a complete JSON object or a single JSON code fence.
It rejects narrative repair, invalid field types, invalid excerpt indices, and truncated outputs.
Failures receive empty answer lists and remain in every denominator.
Normalization applies Unicode NFKC, case folding, and word-token extraction.
It preserves articles and compares deduplicated complete answer sets.
Strict set precision, recall, F1, and exact set accuracy use those normalized sets.
Token overlap remains a secondary measure.

The primary sample contains 60 questions from 55 articles.
Its 600 method conditions reduce to 194 distinct prompts.
The diagnostic retains eleven questions selected by the preceding source version without answer labels or generated answers.
That version selected different possible and guaranteed contexts, up to twenty-four questions.
The revised contexts still differ on seven of those eleven questions.
All eleven remain evaluated, yielding 110 conditions and forty distinct prompts.
They come from nine articles, four shared with the primary sample.
The cohorts have no question overlap and no shared requests.
Of 234 current requests, 166 exactly match archived requests under the same model and execution settings.
The other 68 requests receive new generations.
Reuse checks include full messages, output limits, seed, stop conditions, model hashes, and runtime settings.
Full request hashes, raw responses, timing, and per-request provenance support independent verification.
% Generated only from complete frozen-scoring outputs.
\begin{table}[t]
\centering\footnotesize
\setlength{\tabcolsep}{3pt}
\begin{tabular}{@{}lrrrr@{}}
\toprule
 & \multicolumn{2}{c}{Primary} & \multicolumn{2}{c}{Diagnostic}\\
Method & EM & F1 & EM & F1\\
\midrule
Original BM25 & 16.67 & 16.67 & 9.09 & 9.09\\
Original BM25 + year & 18.33 & 18.33 & 9.09 & 9.09\\
Title-aware ranking & 21.67 & 21.67 & 9.09 & 9.09\\
\quad + earliest completion & 20.00 & 20.00 & 18.18 & 18.18\\
\quad + midpoint completion & 20.00 & 20.00 & 9.09 & 9.09\\
\quad + latest completion & 20.00 & 20.00 & 9.09 & 9.09\\
\quad + interval control & 20.00 & 20.00 & 9.09 & 9.09\\
\quad + possible support & 20.00 & 20.00 & 9.09 & 9.09\\
\quad + guaranteed support & 20.00 & 20.00 & 18.18 & 18.18\\
Title-aware + year & 21.67 & 21.67 & 9.09 & 15.15\\
\bottomrule
\end{tabular}
\caption{Reader exact match (EM) and strict answer-set F1 (\%). Primary: 60 questions and 600 conditions. Diagnostic: 11 questions and 110 conditions. Unique generation parse failures: 69/194 and 11/40, respectively. Failures remain in both denominators.}
\label{tab:reader}
\end{table}


\paragraph{Complete reader outcomes.}
Table~\ref{tab:reader} reports both cohorts separately.
The diagnostic guaranteed-minus-possible F1 difference is 9.09 points, with article-cluster interval $[0.00,30.00]$.
All 45 paired method contrasts remain in each cohort's saved summary.
The primary temporal contrasts change no answer sets, including all 37 pairs with valid responses on both sides.
The diagnostic possible--guaranteed contrast changes three answer sets among seven pairs with two valid responses.
One further change involves a response-format failure.
Earliest--latest completion gives the same counts and changed questions.
Only the Maher change increases strict F1.
The other valid changes concern Atl\'etico Madrid and Helen Clark; both retain zero strict F1.
This diagnostic retains its earlier selection while evaluating the revised source contexts.

Across 234 distinct current requests, eighty responses fail the frozen parser.
These comprise 58 JSON-format failures, four invalid citation lists, and eighteen length-limit responses.
Every failed response remains scored with an empty answer set.
The primary ranking comparison changes twenty answer sets against original BM25 with recency.
Five changes occur among 24 pairs with two valid responses; fifteen involve a parsing failure.
Separate descriptive audits report this breakdown for every method pair without changing the frozen scores.
The release retains 238 earlier executions and 68 new executions, totaling 306 archived executions across both versions.

\subsection{Ranking portability and selection cost}
A secondary analysis applies the certificate to all six ranking variants considered during calibration.
The revised protocol precedes its mechanism analysis and preserves the six variants from the earlier ranking study.
It reads no answer labels and does not replace the selected main ranking.
Each variant has its own modeled-context denominator.
The release also reports fixed main and union cohorts for direct comparisons.
\begin{table*}[t]
\centering\small
\begin{tabular}{lrrrr}
\toprule
Ranking & Template & Human & Template & Human \\
 & stable/modeled & stable/modeled & tests & tests \\
\midrule
Global BM25 & 285/354 & 86/102 & 2.64 & 2.34 \\
Global BM25 + year & 443/534 & 138/164 & 2.72 & 2.39 \\
Title BM25 & 241/298 & 76/86 & 2.66 & 2.41 \\
Title-aware & 189/243 & 56/66 & 2.63 & 2.44 \\
Title BM25 + year & 327/396 & 109/126 & 2.67 & 2.42 \\
Title-aware + year & 281/343 & 102/120 & 2.56 & 2.40 \\
\bottomrule
\end{tabular}
\caption{Secondary certificate analysis across all six prespecified rankings. Each stability denominator contains modeled evidence under that ranking. Tests count temporal predicate calls within that modeled cohort.}
\label{tab:ranking-portability}
\end{table*}

The lazy selector agrees with full-mask evaluation for every parsed query.
Temporal-test counts include possible and guaranteed predicate calls.
Fallback excerpts require no predicate call.
The timing microbenchmarks exclude extraction, compilation, relevance ranking, rendering, and answer generation.
Concurrent CPU workloads affect their absolute timings.
The main text therefore reports operation counts.

\subsection{Research assistance and verification}
AI assistants contributed to code development, theoretical derivations, manuscript drafting, figure construction, and citation checking.
The source interpretation audits are model-assisted judgments.
The release contains independent checks, complete execution records, and a submission-form disclosure covering these uses.

\subsection{Earlier reader and extraction calibration}
Earlier Qwen2.5-1.5B-Instruct runs remain available under their original frozen settings.
A 24-question pilot requires 52 generations for 192 conditions.
A nine-question diagnostic adds 22 generations and reuses six.
The cohorts overlap on two questions, giving 31 distinct questions and 74 actual generations.
All methods reach 31.94\% strict set F1 on the pilot.
The earlier diagnostic changes four answer sets between earliest and latest completion, without improving aggregate strict F1.

Two exploratory learned extractors each process six calibration paragraphs before stopping.
The first returns 41 records, all missing required source quotes.
A schema-constrained variant returns 46 records; 45 fail literal grounding.
The remaining record has an incorrect holder--date association.
Neither extractor supplies claims to the reported confirmation runs.
Partial manifests retain all completed outputs, settings, and stopping decisions.

An earlier deterministic reproduction correction sorts BM25 token accumulation and stops after five accepted units.
Every one of 20,192 archived method contexts remains unchanged.
The original runner and equivalence records are preserved.

\subsection{Complete annual controls}
\begin{table}[t]
\centering
\footnotesize
\setlength{\tabcolsep}{3pt}
\begin{tabular}{lrr}
\toprule
Source-label property & Count & Rate (95\% CI) \\
\midrule
Multiple answers & 6,498/25,752 & 25.2 [24.2, 26.2] \\
Old first retained & 1,456/3,980 & 36.6 [35.1, 38.0] \\
Any answer retained & 6,050/6,961 & 86.9 [86.0, 87.8] \\
\bottomrule
\end{tabular}
\caption{Answer-set properties on held-out source keys. The last two rows condition on first-answer changes and complete-set changes, respectively. Rates and intervals are percentages.}
\label{tab:sourceaudit}
\end{table}

\begin{table}[t]
\centering
\small
\setlength{\tabcolsep}{4pt}
\begin{tabular}{lrrr}
\toprule
Policy & H@1 $\uparrow$ & R@5 $\uparrow$ & S@5 $\downarrow$ \\
\midrule
BM25 & 64.64 & 80.84 & 45.89 \\
BM25 + recent ties & 84.04 & 90.10 & 40.50 \\
Age weighting & 99.81 & 99.98 & 24.07 \\
Date-first ranking & 99.81 & 99.98 & 0.00 \\
Original predicate & 99.37 & 99.77 & 0.77 \\
Graphiti kernel & 99.78 & 99.93 & 0.00 \\
Query latest-value & 99.89 & 86.37 & 0.00 \\
Query latest-batch & 99.89 & 99.98 & 0.00 \\
\bottomrule
\end{tabular}
\caption{Annual source-label retrieval on 25,752 held-out questions. All values are percentages. H@1: first-result accuracy; R@5: answer recall; S@5: target-label stale exposure.}
\label{tab:fullsource}
\end{table}

The annual split and complete-source labels match the archived experiment.
Latest-value selection retains 99.89\% first-result accuracy but achieves 86.37\% complete answer recall.
Latest-batch selection reaches 99.98\% recall.
The suppression difference is 13.63 points, with key-cluster interval $[13.06,14.21]$.
Graphiti's matched policy kernel and direct-first invalidation produce identical endpoints across 45,866 observations.
These retained results require no repeated model inference.

\subsection{Earlier natural paragraph controls}
\begin{table}[t]
\centering
\small
\setlength{\tabcolsep}{3.5pt}
\begin{tabular}{lrrrr}
\toprule
& \multicolumn{2}{c}{Human} & \multicolumn{2}{c}{Templates} \\
Policy & H@5 & R@5 & H@5 & R@5 \\
\midrule
BM25 & 45.66 & 45.04 & 32.49 & 31.75 \\
Latest year & 54.70 & 54.02 & 51.86 & 50.89 \\
Year envelope & 49.16 & 48.33 & 34.44 & 33.67 \\
Explicit range & 44.10 & 43.53 & 32.11 & 31.39 \\
Range support & 44.58 & 44.00 & 38.50 & 37.63 \\
\bottomrule
\end{tabular}
\caption{TimeQA paragraph retrieval percentages. H@5 measures annotated-passage coverage; R@5 measures annotated-span recall. The tracks contain 830 human questions and 2,613 template questions.}
\label{tab:natural}
\end{table}

These paragraph runs use the published test passages and positive annotations.
Their human-question results were inspected before the source-unit experiments.
Paragraph and source-unit runs have different retrieval units and denominators.
